\documentclass[10pt]{article}
\usepackage[T1]{fontenc}
\usepackage[utf8]{inputenc}
\usepackage[french]{babel}
\usepackage{amsmath,amssymb,mathrsfs,mathpazo}
\usepackage{geometry}
\geometry{a4paper,margin=2cm}
\usepackage{graphicx,array,tabularx,tabularray,colortbl,multirow,multicol,booktabs,enumitem,pifont,ulem,fancybox,diagbox,cancel}
\usepackage{pstricks,pstricks-add,pst-plot,pst-tree,pst-node,pst-text,pst-func,pst-math,pst-bspline,pst-3dplot,pst-eucl,pst-solides3d}
\let\pscalculate\undefined
\usepackage{pst-bezier}
\usepackage{tikz,pgfplots,tkz-tab}
\usetikzlibrary{arrows,arrows.meta,patterns,calc,positioning,trees,decorations.markings,shapes.geometric,shapes.misc,shapes.arrows,intersections,decorations.pathmorphing,decorations.pathreplacing,backgrounds,shadows,babel}
\pgfplotsset{compat=1.18}
\usepgfplotslibrary{fillbetween}
\usepackage[np]{numprint}
\usepackage{listings,makecell,pas-tableur,twoopt}
\let\barre\undefined
\usepackage[pstricks]{tablvar}
\definecolor{darkgreen}{RGB}{0,100,0}
\providecommand{\red}{\color{red}}
\providecommand{\blue}{\color{blue}}
\providecommand{\green}{\color{green}}
\providecommand{\black}{\color{black}}
\providecommand{\white}{\color{white}}
\providecommand{\gray}{\color{gray}}
\providecommand{\cyan}{\color{cyan}}
\providecommand{\magenta}{\color{magenta}}
\providecommand{\yellow}{\color{yellow}}
\providecommand{\violet}{\color{violet}}
\providecommand{\orange}{\color{orange}}
\providecommand{\euro}{\texteuro}
\providecommand{\EUR}{\texteuro}
\providecommand{\ieme}{\textsuperscript{e}}
\providecommand{\degres}{\ensuremath{^\circ}}
\providecommand{\degree}{\ensuremath{^\circ}}
\providecommand{\vv}[1]{\overrightarrow{#1}}
\providecommand{\Coord}[1]{#1}
\providecommand{\decosix}{}
\providecommand{\starredbullet}{\ensuremath{\bullet}}
\pagestyle{empty}
\setlength{\parindent}{0pt}
\usepackage{adjustbox}\DeclareUnicodeCharacter{2236}{:}\DeclareUnicodeCharacter{2190}{\ensuremath{\leftarrow}}
\providecommand{\R}{\mathbb{R}}
\providecommand{\N}{\mathbb{N}}
\providecommand{\Z}{\mathbb{Z}}
\providecommand{\Q}{\mathbb{Q}}
\providecommand{\D}{\mathbb{D}}
\providecommand{\C}{\mathbb{C}}
\providecommand{\e}{{\mathrm{e}}}
\providecommand{\vect}[1]{\overrightarrow{#1}}
\providecommand{\vectt}[1]{\overrightarrow{\mathrm{#1}}}
\providecommand{\vvt}[1]{\overrightarrow{\mathrm{#1}}}
\providecommand{\np}[2][]{{#2~#1}}
\providecommand{\vv}[1]{{\overrightarrow{#1}}}
\renewcommand{\d}{{\mathrm{d}}}
\providecommand{\degres}{{^\circ}}
\providecommand{\degre}{{^\circ}}
\providecommand{\blue}{\color{blue}}
\providecommand{\red}{\color{red}}
\providecommand{\enskip}{\quad}
\providecommand{\llbracket}{\lbrack\!\lbrack}
\providecommand{\rrbracket}{\rbrack\!\rbrack}
\providecommand{\euro}{{\text{€}}}
\providecommand{\no}{{\text{n°}}}
\providecommand{\interval}[2]{{[#1;#2]}}
\providecommand{\intervFO}[2]{{[#1;#2[}}
\providecommand{\dotfill}{\ldots}

\providecommand{\ds}{\displaystyle}
\providecommand{\textding}[1]{#1}
\providecommand{\Oij}{\(\mathrm{O};\vec{\imath},\vec{\jmath}\)}
\providecommand{\Oijk}{\(\mathrm{O};\vec{\imath},\vec{\jmath},\vec{k}\)}
\providecommand{\Ouv}{\(\mathrm{O};\vec{u},\vec{v}\)}
\providecommand{\intervFF}[2]{\ensuremath{[#1;#2]}}
\renewcommand{\euro}{\ifmmode\text{\texteuro}\else\texteuro\fi}
\DeclareUnicodeCharacter{00B0}{\ifmmode^{\circ}\else\textdegree\fi}
\DeclareUnicodeCharacter{20AC}{\ifmmode\text{\texteuro}\else\texteuro\fi}
% Énoncés : Éducation nationale. Transcription et corrigés : APMEP, Denis Vergès et contributeurs. Découpage et classement : CoopMaths, Sébastien Lozano, Éric Elter et contributeurs. Adaptation de la présentation : Maths Explorer.
% CC BY-SA : https://coopmaths.fr/www/about/
% https://coopmaths.fr/alea/static/bac/2022/tex/bac_2022_09_sujet2_metropole_1_cor.tex
\begin{document}
\noindent\textit{APMEP / CoopMaths. CC BY-SA. Présentation adaptée par Maths Explorer.}\par\medskip
\medskip

Dans le magasin d'Hugo, les clients peuvent louer deux types de vélos : vélos de route ou bien vélos tout terrain. 

Chaque type de vélo peut être loué dans sa version électrique ou non.

On choisit un client du magasin au hasard, et on admet que:
\begin{itemize}
\item Si le client loue un vélo de route, la probabilité que ce soit un vélo électrique est de 0,4 ;
\item Si le client loue un vélo tout terrain, la probabilité que ce soit un vélo électrique est de 0,7 ; 
\item La probabilité que le client loue un vélo électrique est de $0,58$.
\end{itemize}

\medskip

On appelle $\alpha$ la probabilité que le client loue un vélo de route, avec $0 \leqslant \alpha \leqslant 1$. 

On considère les évènements suivants:

\begin{itemize}
\item[$\bullet~~$] $R$ : \og le client loue un vélo de route \fg{} ;
\item[$\bullet~~$] $E$ : \og le client loue un vélo électrique \fg{} ;
\item[$\bullet~~$] $\overline{R}$ et $\overline{E}$, évènements contraires de $R$ et $E$.
\end{itemize}

























\begin{enumerate}
\item On complète l'arbre proposé.

\begin{center}
\pstree[treemode=R,nodesepA=0pt,nodesepB=2.5pt,treesep = 1cm, levelsep=3cm,nrot=:U]{\TR{}}
{\pstree{\TR{$R$~}\naput{$\alpha$}}
{\TR{$E$}\naput{\blue $0,4$}
\TR{$\overline{E}$}\nbput{\blue $1-0,4=0,6$}
}
\pstree{\TR{$\overline{R}$~}\nbput{$1 - \alpha$}}
{\TR{$E$}\naput{\blue $0,7$}
\TR{$\overline{E}$}\nbput{\blue$1-0,7=0,3$}
}}
\bigskip
\end{center}

\item 
	\begin{enumerate}
		\item D'après la formule des probabilités totales:
		
		 $p(E) = p\left (R\cap E\right )+p\left (\overline{R}\cap E\right )
		 = \alpha \times 0,4+\left (1-\alpha\right )\times 0,7
		 =0,4\alpha + 0,7 - 0,7\alpha 
		 = 0,7 - 0,3\alpha$.
		 
		\item  La probabilité que le client loue un vélo électrique est  $p(E) = 0,58$.
		
Or $p(E)=0,7-0,3\alpha$. Donc $0,7 - 0,3\alpha = 0,58$ ce qui équivaut à $0,7-0,58 = 0,3\alpha$ ou encore $0,12=0,3\alpha$ soit $\alpha = 0,4$.
	\end{enumerate}
	
\item On sait que le client a loué un vélo électrique. 

La probabilité qu'il ait loué un vélo tout terrain est:

$p_E\left (\overline{R}\right ) 
= \dfrac{p\left ( \overline{R}\cap E\right )}{p(E)}
= \dfrac{\left (1-0,4\right )\times 0,7}{0,58}
= \dfrac{0,42}{0,58} \approx 0,72$.



\item La probabilité que le client loue un vélo tout terrain électrique est:
$p\left (\overline{R} \cap E\right )=0,42$.

\item Le prix de la location à la journée d'un vélo de route non électrique est de $25$~euros, celui d'un vélo tout terrain non électrique de $35$~euros. 
Pour chaque type de vélo, le choix de la version électrique augmente le prix de location à la journée de $15$~euros. 

On appelle $X$ la variable aléatoire modélisant le prix de location d'un vélo à la journée. 
	\begin{enumerate}
		\item 
On a quatre possibilités.

\begin{list}{\textbullet}{}
\item La location d'un vélo de route non électrique coûte 25~\euro.

Cela correspond à  l'évènement $R\cap \overline{E}$ de probabilité $0,4\times 0,6=0,24$.

\item La location d'un vélo de route  électrique coûte $25+15$ soit 40~\euro.

Cela correspond à  l'évènement $R\cap E$ de probabilité $0,4\times 0,4=0,16$.
 
\item La location d'un vélo tout terrain non électrique coûte 35~\euro.

Cela correspond à  l'évènement $\overline{R} \cap \overline{E}$ de probabilité $0,6 \times 0,3=0,18$.

\item La location d'un vélo tout terrain  électrique coûte $35+15$ soit 50~\euro.

Cela correspond à  l'évènement $\overline{R}\cap E$ de probabilité $0,6\times 0,7=0,42$.
 
\end{list}		
		
On établit la loi de probabilité de $X$:

\begin{center}
\begin{tabularx}{0.7\linewidth}{|c|*{4}{>{\centering \arraybackslash}X|}}
\hline
$x_i$ & 25 & 35 & 40 & 50\\
\hline
$p_i=p(X=x_i)$ & $0,24$ & $0,18$ & $0,16$ & $0,42$\\
 \hline
\end{tabularx}
\end{center}
			
		\item 
L'espérance mathématique de $X$ est:

$E(X)=\ds\sum x_i\times p_i 
= 25\times 0,24 + 35\times 0,18 + 40\times 0,16 + 50\times 0,42 
= 39,7$.		

Le coût moyen d'une location est donc de $39,70$~euros.		
		
	\end{enumerate}	
\item Lorsqu'on choisit $30$ clients d'Hugo au hasard, on assimile ce choix à un tirage avec remise. 
On note $Y$ la variable aléatoire associant à un échantillon de 30 clients choisis au hasard le nombre de clients qui louent un vélo électrique.

On rappelle que la probabilité de l'évènement $E$ est : $p(E) = 0,58$.
	\begin{enumerate}
		\item 
Il s'agit d'une répétition de 30 épreuves identiques et indépendantes n'ayant que deux issues, la probabilité du succès pour une épreuve étant égale à $0,58$.		
		
		Donc la variable aléatoire $Y$ qui donne le nombre de succès sur 30 tirages, suit une loi binomiale de paramètres $n=30$ et $p=0,58$.
				
		\item La probabilité qu'un échantillon contienne exactement $20$ clients qui louent un vélo électrique est:
		
$p(Y=20) = \ds\binom{30}{20}\times 0,58^{20}	\times (1-0,58)^{30-20}
\approx 0,095$.

		\item La probabilité qu'un échantillon contienne au moins $15$ clients qui louent un vélo électrique est:

$p(Y\geqslant 15) = 1 - p(Y\leqslant 14)\approx 1 - 0,14190 \approx 0,858$.

	\end{enumerate}

\end{enumerate}

\end{document}
